% !TEX program = xelatex
% ============================================================
%  Quran & Hadith Template
%  Arabic (RTL) LaTeX template for Islamic documents with
%  Quran verses, Hadith citations, and scholarly commentary.
%
%  Requires: XeLaTeX, Amiri + NotoKufiArabic-Regular.ttf fonts, quran package
%  Compile:  xelatex main.tex && xelatex main.tex
% ============================================================
\documentclass[12pt,a4paper]{article}

% --- Page geometry ---
\usepackage[a4paper,margin=2.5cm,top=2.5cm,bottom=2.5cm,headheight=15pt]{geometry}

% --- Standard packages (load ALL before bidi) ---
\usepackage{array}
\usepackage{booktabs}
\usepackage{tabularx}
\usepackage{enumitem}
\usepackage{float}
\usepackage{titlesec}
\usepackage{fancyhdr}
\usepackage{setspace}
\usepackage[table]{xcolor}

% --- BIDI (must come before polyglossia, after standard packages) ---
\usepackage{bidi}

% --- Font specification (requires bidi) ---
\usepackage[no-math]{fontspec}

% --- Polyglossia (language support) ---
\usepackage[quiet]{polyglossia}

% --- Fonts ---
\setmainfont[Scale=1]{NotoKufiArabic-Regular.ttf}
\newfontfamily\arabicfont[Script=Arabic,Scale=0.95]{NotoKufiArabic-Regular.ttf}
\newfontfamily\englishfont[Scale=0.95]{TeX Gyre Termes}
% Specialized font for Quranic text (better diacritics, line height)
\newfontfamily\quranfont[Script=Arabic,Scale=1.15]{AmiriQuran.ttf}

% --- Languages ---
\setdefaultlanguage[calendar=gregorian,hijricorrection=1]{arabic}
\setotherlanguage[variant=british]{english}

% --- Quran package ---
% Options: uthmani (full diacritics, may exceed memory on some systems)
%          uthmani-min (minimal uthmani)
%          nopar (no paragraph breaks between ayahs)
%          nonumber (hide ayah numbers)
%          wordwise (enable chunk-level typesetting)
%          ornbraces (use ornamental braces ﴿ ﴾)
%          basmalah / nobasmalah (control basmalah display)
\usepackage[ornbraces]{quran}

% --- Hyperref (ALWAYS last) ---
\usepackage[breaklinks,hidelinks]{hyperref}
\hypersetup{pdftitle={القرآن والحديث},pdfauthor={المؤلف}}

% --- Section formatting ---
\titleformat{\section}{\Large\bfseries}{\thesection}{0.7em}{}
\titleformat{\subsection}{\large\bfseries}{\thesubsection}{0.6em}{}
\titlespacing*{\section}{0pt}{1.5ex plus 0.3ex}{0.8ex plus 0.2ex}

% --- Header/Footer ---
\pagestyle{fancy}
\fancyhf{}
\fancyhead[R]{\small\itshape القرآن والحديث}
\fancyhead[L]{\small اسم المؤلف}
\fancyfoot[C]{\small\thepage}
\renewcommand{\headrulewidth}{0.3pt}

% --- Table column types ---
\newcolumntype{L}[1]{>{\raggedright\arraybackslash}p{#1}}
\newcolumntype{C}[1]{>{\centering\arraybackslash}p{#1}}
\newcolumntype{R}[1]{>{\raggedleft\arraybackslash}p{#1}}

% --- Line spacing ---
\setstretch{1.6}

% ============================================================
%  Custom Commands for Islamic Documents
% ============================================================

% --- Quran verse with decorative formatting ---
% Usage: \quranverse{surah_number}{ayah_number}
% Renders a single ayah in the quranfont with ornamental ending
\newcommand{\quranverse}[2]{%
  \begin{quote}
    {\quranfont\setRTL\quranayah[#1][#2]}%
  \end{quote}
}

% --- Quran surah (full surah) ---
% Usage: \quransurahdisplay{surah_number}
\newcommand{\quransurahdisplay}[1]{%
  \begin{quote}
    {\quranfont\setRTL\quransurah*[#1]}%
  \end{quote}
}

% --- Hadith with isnad and source ---
% Usage: \hadith{Arabic text}{Source/grade}
\newcommand{\hadith}[2]{%
  \begin{quote}
    \textarabic{\large #1}\par
    \vspace{0.3em}
    \hfill\textit{\small #2}
  \end{quote}
}

% --- Hadith with isnad, matn, and source ---
% Usage: \hadithfull{Isnad}{Matn}{Source}
\newcommand{\hadithfull}[3]{%
  \begin{quote}
    \textarabic{\large #1}\par
    \vspace{0.2em}
    \textarabic{\large #2}\par
    \vspace{0.3em}
    \hfill\textit{\small #3}
  \end{quote}
}

% --- Quran citation (inline, not footnote to avoid bidi/hyperref memory issues) ---
% Usage: \qurancite{surah_name_or_number}{ayah_number}
\newcommand{\qurancite}[2]{\textit{[القرآن الكريم، سورة #1، آية #2]}}

% --- Hadith citation (inline) ---
% Usage: \hadithcite{collection}{number}
\newcommand{\hadithcite}[2]{\textit{[#1، حديث #2]}}

% --- Basmalah (opening phrase) ---
\newcommand{\basmalah}{%
  \begin{center}
    {\quranfont\setRTL بِسْمِ اللَّهِ الرَّحْمَٰنِ الرَّحِيمِ}
  \end{center}
}

% --- Decorative separator ---
\newcommand{\separator}{%
  \begin{center}
    \rule{0.3\textwidth}{0.5pt}
  \end{center}
}

% ============================================================
%  Document
% ============================================================
\begin{document}

% --- Title ---
\begin{center}
{\LARGE\bfseries القرآن الكريم والحديث الشريف}\\[0.5em]
{\large قالب للوثائق الإسلامية}\\[1em]
\separator
\end{center}

\vspace{1em}

% ============================================================
\section{المقدمة}
% ============================================================

هذا قالبٌ لإعداد الوثائق الإسلامية التي تحتوي على آيات قرآنية وأحاديث نبوية مع التعليق العلمي. يستخدم القالب حزمة \LR{\texttt{quran}} لاستدعاء الآيات مباشرةً، وخط \LR{NotoKufiArabic-Regular.ttf} لعرض النص القرآني بشكلٍ مناسب.

\separator

% ============================================================
\section{أمثلة من القرآن الكريم}
% ============================================================

\subsection{سورة الفاتحة}

\basmalah

\quransurahdisplay{1}

\separator

\subsection{آية الكرسي}

\quranverse{2}{255}

قال تعالى: \LR{``}\textarabic{اللَّهُ لَا إِلَٰهَ إِلَّا هُوَ الْحَيُّ الْقَيُّومُ}\LR{''}\qurancite{البقرة}{255}

\separator

\subsection{سورة الإخلاص}

\quransurahdisplay{112}

\separator

\subsection{آيات مختارة}

\quranverse{1}{1}

\quranverse{36}{1}

\quranverse{55}{1}

\separator

% ============================================================
\section{أمثلة من الحديث الشريف}
% ============================================================

\subsection{حديث النية}

\hadith{إِنَّمَا الْأَعْمَالُ بِالنِّيَّاتِ، وَإِنَّمَا لِكُلِّ امْرِئٍ مَا نَوَى}{متفق عليه، رواه البخاري ومسلم}

\separator

\subsection{حديث الإيمان}

\hadith{مَنْ كَانَ يُؤْمِنُ بِاللَّهِ وَالْيَوْمِ الْآخِرِ فَلْيَقُلْ خَيْرًا أَوْ لِيَصْمُتْ}{متفق عليه}

\separator

\subsection{حديث مع الإسناد}

\hadithfull{%
  عَنْ عُمَرَ بْنِ الْخَطَّابِ رَضِيَ اللَّهُ عَنْهُ قَالَ:%
}{%
  بَيْنَمَا نَحْنُ جُلُوسٌ عِنْدَ رَسُولِ اللَّهِ صَلَّى اللَّهُ عَلَيْهِ وَسَلَّمَ ذَاتَ يَوْمٍ...%
}{%
  أخرجه الإمام مسلم في صحيحه، حديث رقم 8%
}

\separator

% ============================================================
\section{جدول المراجع}
% ============================================================

\begin{table}[H]
\centering
\begin{tabularx}{\textwidth}{L{5cm} C{3cm} X}
\toprule
\textbf{المصدر} & \textbf{الرقم} & \textbf{الوصف} \\
\midrule
القرآن الكريم & --- & سورة الفاتحة، البقرة، الإخلاص \\
صحيح البخاري & 1 & كتاب بدء الوحي \\
صحيح مسلم & 8 & كتاب الإيمان \\
\bottomrule
\end{tabularx}
\caption{جدول المراجع الإسلامية}
\end{table}

% ============================================================
\section{ملاحظات الاستخدام}
% ============================================================

\subsection{استدعاء الآيات}

استخدم الأوامر التالية لاستدعاء الآيات:

\begin{itemize}[itemsep=0.3em]
  \item \LR{\texttt{\textbackslash quranverse\{2\}\{255\}}} --- آية واحدة (السورة 2، الآية 255)
  \item \LR{\texttt{\textbackslash quransurahdisplay\{1\}}} --- سورة كاملة (الفاتحة)
  \item \LR{\texttt{\textbackslash quranayah[2][255]}} --- الأمر الأصلي من حزمة \LR{quran}
  \item \LR{\texttt{\textbackslash quransurah[112]}} --- سورة كاملة كفقرة واحدة
\end{itemize}

\subsection{استدعاء الأحاديث}

استخدم الأوامر التالية لإدراج الأحاديث:

\begin{itemize}[itemsep=0.3em]
  \item \LR{\texttt{\textbackslash hadith\{النص\}\{المصدر\}}} --- حديث مع المصدر
  \item \LR{\texttt{\textbackslash hadithfull\{الإسناد\}\{المتن\}\{المصدر\}}} --- حديث كامل بالإسناد
  \item \LR{\texttt{\textbackslash hadithcite\{صحيح البخاري\}\{1\}}} --- اقتباس في الحاشية
\end{itemize}

\subsection{خيارات حزمة Quran}

يمكن تعديل خيارات حزمة \LR{quran} في السطر:

\begin{quote}
\LR{\texttt{\textbackslash usepackage[ornbraces]\{quran\}}}
\end{quote}

الخيارات المتاحة:

\begin{itemize}[itemsep=0.3em]
  \item \LR{\texttt{uthmani}} --- نص عثماني كامل بالتشكيل (قد يستهلك ذاكرة كبيرة)
  \item \LR{\texttt{uthmani-min}} --- نص عثماني مبسط
  \item \LR{\texttt{nopar}} --- بدون فواصل فقرات بين الآيات
  \item \LR{\texttt{nonumber}} --- إخفاء أرقام الآيات
  \item \LR{\texttt{wordwise}} --- تمكين استدعاء أجزاء من الآية
  \item \LR{\texttt{ornbraces}} --- استخدام الأقواس الزخرفية ﴿ ﴾
  \item \LR{\texttt{basmalah}} / \LR{\texttt{nobasmalah}} --- التحكم في عرض البسملة
\end{itemize}

\subsection{الترجمات}

يمكن إضافة ترجمات القرآن باستخدام الحزم الإضافية:

\begin{itemize}[itemsep=0.3em]
  \item \LR{\texttt{\textbackslash usepackage[transen]\{quran\}}} --- ترجمة إنجليزية
  \item \LR{\texttt{\textbackslash usepackage\{quran-en\}}} --- 15 ترجمة إنجليزية إضافية
  \item \LR{\texttt{\textbackslash usepackage\{quran-ur\}}} --- ترجمات أردية
  \item \LR{\texttt{\textbackslash usepackage\{quran-de\}}} --- ترجمات ألمانية
\end{itemize}

\separator

\begin{center}
\textit{تم بحمد الله}
\end{center}

\end{document}
